\section{Frontier Coding, Bounded}
\sectionkicker{FRONTIER CODING}
\begin{wideflow}
{\Large\bfseries コードの旗艦は限定付きで受け止める}\par
\smallskip
{\color{SurveyMuted}三つのコード系の出荷を、手法の非開示とともに並べる。}\par
\end{wideflow}

9月にAnthropicが示したClaude Fable 5.1とClaude Mythos 5.1は、同じモデルをセーフガード水準の違いで出す位置づけである\cite{fablemythos}。Fable 5.1は一般提供、Mythos 5.1はサイバーと生命科学のタスク向けの信頼されたアクセスプログラム経由とされる。Fable 5.1の価格は、キャッシュ読み取りが100万トークンあたり0.25ドル（75\%減）、通常のトークン課金タスクで約25\%安く、エージェント性の高いタスクで最大約45\%安いとされ、それ以外は入力10ドル、出力50ドル、API名claude-fable-5-1である。測定値としてTerminal-Bench-Science 52.6\%、Terminal-Bench 4.0 55.8\%（Mythosは60.9\%）、GDPval-AA v2 1853、OSWorld 2.0部分一致77.9\%・厳密一致41.7\%、AutomationBench 31.4\%、CursorBench 73.4\%が掲げられる。安全対策では、Mythosの化学・生物の能力がMythos 5を上回り次のRSPリスク段階には届かず、最も強いAnthropicのサイバー能力でありながら下位のFCFリスク区分で重大な脱獄は見つからず、生物分野のセーフガードが良性クエリに反応する割合が85\%減、サイバー分野の安全対策は平均で約6割の介入削減とされる。いずれも研究室の報告であり、外部の再現はない。Mythosの審査基準や透かし検出APIの試用範囲は発表されたプログラムであり、稼働中の配備の観察ではない。科学分野の事例紹介は提供側が選んだ報告であり、統制された証拠ではない。

\begin{claimboundary}[報告の境界]
測定値と安全対策の数値は研究室の報告であり、外部の再現を待つ。審査基準や透かし検出は発表されたプログラムであり、配備の観察ではない。
\end{claimboundary}

9月2日にGoogleが出したGemini 3.8 FlashとGemini 3.8 Flash Cyberは、同じ基礎能力を二つの形態で出す位置づけである\cite{gemini38flash}。Flashは汎用の主力、Flash CyberはFairwind Program経由で信頼された防御者向けとされる。Flashの導入価格は2026年12月31日まで入力0.75ドル、出力3.75ドル（100万トークンあたり）、2027年1月1日から入力1.50ドル、出力7.50ドルである。能力主張は、長い期間のコード生成やエージェント利用の伸び、Vals Finance Agent V2やHarvey Legalでの3.7 Flashや他の旗艦モデル超え、HLE-Verified 54.9\%、高いエフォート段階での追加の推論と丁寧なツール呼び出し、トークン増加とされる。Flash Cyberは、CyberGymでの最先端級の発見、内部の20言語での脆弱性発見で7割超、CWE-Bench pass@1で47.2\%対最先端47.8\%（低コスト）、Chromeの2.6倍の正しいパッチ、Wizのリコール7.5から9.7ポイント増（2.3から5.2倍の低コスト）、2時間以内でのGoogleの重大な発見1件とされる。測定値やパッチのリコール、発見の数値は手法の確認なしの提供側報告であり、独立の再現を待つ。内部の20言語ベンチマークは非開示であり、CyberGymやCWE-Benchの実行はGoogleの報告である。WizやChromeチームの帰属はGoogle経由の伝聞であり、独立に確かめていない。DeepMindのモデルカードの追加詳細は本号の事実確認には用いない。

9月2日にMetaが出したMuse Spark 1.3は、Muse CodeとMeta Model API（dev.meta.ai）上で最大推論設定で使えるとされる\cite{musespark}。提供側のコード・ツール利用の伸びとして、確認質問と影響の大きい動作の確認を伴う長い協働、長い指示の保持、シングルスレッドでの複数処理、能力の自己認識、Muse Spark 1.2比でツール呼び出し約2割減・トークン約25\%減という整ったコードが挙げられる。安全対策の伸びとして、敵対的入力やプロンプト注入への耐性と、複雑なエージェント利用タスクでの不可逆動作の調整が挙げられる。いずれも提供側の報告であり、再現はない。ページに埋め込まれた測定値の図表の数値は取得できたmarkdownからは読み取れず、関連するマルチモーダル評価の手法報告も確認していないため、測定値は本号の本文に取り込まない。重みの公開や大型化は将来構想の言葉であり、出荷の事実ではない。

\begin{technicalnote}[図表数値の扱い]
Museの図表の数値と手法報告は確認していないため、測定値は取り込まない。将来構想の言葉を出荷の事実として扱わない\cite{musespark}。
\end{technicalnote}
